\documentclass[11pt,a4paper]{article}
\usepackage[utf8]{inputenc}
\usepackage[T1]{fontenc}
\usepackage{lmodern}
\usepackage[french]{babel}
\usepackage{amsmath,amssymb,amsthm}
\usepackage[margin=2.4cm]{geometry}
\usepackage{xcolor}
\usepackage{tcolorbox}
\usepackage{tikz}
\usepackage{pgfplots}
\pgfplotsset{compat=1.18}
\usetikzlibrary{arrows.meta}

\setlength{\parindent}{0pt}
\setlength{\parskip}{4pt}
\numberwithin{equation}{section}

\theoremstyle{plain}
\newtheorem{proposition}{Proposition}[section]
\newtheorem{lemme}[proposition]{Lemme}
\theoremstyle{definition}
\newtheorem{definition}[proposition]{Définition}

\newcommand{\sign}{\operatorname{sign}}
\newcommand{\hth}{\hat{\theta}}

\newtcolorbox{idee}{colback=blue!4,colframe=blue!50!black,title=Idée,fonttitle=\bfseries}
\newtcolorbox{resultat}{colback=green!5,colframe=green!40!black,title=Résultat,fonttitle=\bfseries}
\newtcolorbox{attention}{colback=red!4,colframe=red!50!black,title=Attention,fonttitle=\bfseries}

\title{\textbf{Démonstration détaillée de}\\[4pt]
$\sign(\varepsilon)=\sign(\psi)$ \textbf{pour} $0<|\psi|<\pi$}
\date{}

\begin{document}
\maketitle

%=====================================================================
\section{Énoncé et notations}
%=====================================================================

On rappelle que le détecteur de phase donne
\begin{equation}\label{eq:eps}
\varepsilon = V_m\sin\psi,\qquad \psi = \theta-\hth,\qquad V_m>0 .
\end{equation}

\begin{proposition}\label{prop:main}
Pour tout $\psi$ tel que $0<|\psi|<\pi$ :
\begin{equation}
\sign(\varepsilon) = \sign(\psi).
\end{equation}
\end{proposition}

\begin{idee}
La démonstration se fait en trois étapes :
\begin{enumerate}
  \item on montre que $V_m$ n'influence pas le signe : $\sign(\varepsilon)=\sign(\sin\psi)$ ;
  \item on montre que $\sin\psi>0$ quand $0<\psi<\pi$ ;
  \item on en déduit que $\sin\psi<0$ quand $-\pi<\psi<0$, grâce à l'imparité du sinus.
\end{enumerate}
\end{idee}

%=====================================================================
\section{Rappels sur la fonction signe}
%=====================================================================

\begin{definition}
\begin{equation}
\sign(x) = \begin{cases} +1 & \text{si } x>0,\\ \phantom{+}0 & \text{si } x=0,\\ -1 & \text{si } x<0. \end{cases}
\end{equation}
\end{definition}

\begin{lemme}[Signe d'un produit]\label{lem:produit}
Pour tous réels $a$ et $b$ : $\sign(ab) = \sign(a)\,\sign(b)$.
\end{lemme}
\begin{proof}
On examine tous les cas possibles.
\begin{center}
\renewcommand{\arraystretch}{1.25}
\begin{tabular}{|c|c|c|c|c|}
\hline
$a$ & $b$ & $ab$ & $\sign(ab)$ & $\sign(a)\sign(b)$\\
\hline
$>0$ & $>0$ & $>0$ & $+1$ & $(+1)(+1)=+1$\\
$>0$ & $<0$ & $<0$ & $-1$ & $(+1)(-1)=-1$\\
$<0$ & $>0$ & $<0$ & $-1$ & $(-1)(+1)=-1$\\
$<0$ & $<0$ & $>0$ & $+1$ & $(-1)(-1)=+1$\\
\multicolumn{2}{|c|}{$a=0$ ou $b=0$} & $=0$ & $0$ & $0$\\
\hline
\end{tabular}
\end{center}
Les deux dernières colonnes sont identiques dans tous les cas.
\end{proof}

\begin{lemme}[Multiplication par une constante positive]\label{lem:const}
Si $c>0$, alors $\sign(c\,x) = \sign(x)$.
\end{lemme}
\begin{proof}
D'après le lemme~\ref{lem:produit}, $\sign(cx) = \sign(c)\sign(x)$. Or $c>0$, donc $\sign(c)=+1$
et $\sign(cx) = \sign(x)$.
\end{proof}

%=====================================================================
\section{Étape 1 : l'amplitude $V_m$ ne change pas le signe}
%=====================================================================

\textbf{Étape 1.a.} On part de \eqref{eq:eps} :
\begin{equation}
\sign(\varepsilon) = \sign(V_m\sin\psi).
\end{equation}
\textbf{Étape 1.b.} On applique le lemme~\ref{lem:const} avec $c=V_m>0$ et $x=\sin\psi$ :
\begin{equation}\label{eq:etape1}
\boxed{\sign(\varepsilon) = \sign(\sin\psi).}
\end{equation}

\begin{attention}
Cette étape utilise $V_m>0$. Si la tension du réseau s'annule ($V_m=0$, par exemple lors d'un court-circuit
franc), alors $\varepsilon=0$ quel que soit $\psi$ : le détecteur ne donne plus aucune information.
\end{attention}

Il reste donc à montrer que $\sign(\sin\psi) = \sign(\psi)$ pour $0<|\psi|<\pi$.

%=====================================================================
\section{Étape 2 : $\sin\psi>0$ pour $0<\psi<\pi$}
%=====================================================================

On donne deux démonstrations : une géométrique (plus intuitive), une analytique (plus rigoureuse).

\subsection{Démonstration géométrique (cercle trigonométrique)}

Sur le cercle de rayon 1, le point d'angle $\psi$ a pour coordonnées $(\cos\psi,\sin\psi)$.
\textbf{$\sin\psi$ est donc la hauteur (l'ordonnée) du point.}

\begin{center}
\begin{tikzpicture}[scale=2.3,>=Stealth]
  \fill[green!12] (0,0) -- (1,0) arc (0:180:1) -- cycle;
  \fill[red!10]   (0,0) -- (1,0) arc (0:-180:1) -- cycle;
  \draw[->] (-1.35,0) -- (1.4,0) node[right] {$\cos$};
  \draw[->] (0,-1.3) -- (0,1.35) node[above] {$\sin$};
  \draw[thick] (0,0) circle (1);
  % point dans le demi-plan supérieur
  \coordinate (P) at ({cos(55)},{sin(55)});
  \draw[very thick,blue,->] (0,0) -- (P);
  \fill[blue] (P) circle (0.03);
  \draw[dashed,blue] (P) -- ({cos(55)},0);
  \draw[blue,line width=2pt] ({cos(55)},0) -- (P);
  \node[blue,right] at ({cos(55)},{sin(55)/2}) {$\sin\psi>0$};
  \draw[->,blue] (0.3,0) arc (0:55:0.3);
  \node[blue] at (0.42,0.18) {$\psi$};
  % point dans le demi-plan inférieur
  \coordinate (Q) at ({cos(-125)},{sin(-125)});
  \draw[very thick,red!70!black,->] (0,0) -- (Q);
  \fill[red!70!black] (Q) circle (0.03);
  \draw[red!70!black,line width=2pt] ({cos(-125)},0) -- (Q);
  \node[red!70!black,left] at ({cos(-125)},{sin(-125)/2}) {$\sin\psi<0$};
  \draw[->,red!70!black] (0.22,0) arc (0:-125:0.22);
  \node[red!70!black] at (0.05,-0.33) {$\psi$};
  % étiquettes
  \node[green!40!black,align=center] at (-1.45,0.9) {\small $0<\psi<\pi$\\ \small (demi-cercle du haut)};
  \node[red!60!black,align=center] at (1.5,-0.9) {\small $-\pi<\psi<0$\\ \small (demi-cercle du bas)};
  \fill (1,0) circle (0.025) node[below right] {\small $\psi=0$};
  \fill (-1,0) circle (0.025) node[below left] {\small $\psi=\pm\pi$};
\end{tikzpicture}
\end{center}

\begin{itemize}
  \item Quand $\psi$ va de $0$ à $\pi$ (sens trigonométrique), le point parcourt le \textbf{demi-cercle du haut}.
        Il est strictement au-dessus de l'axe horizontal, donc $\sin\psi>0$.
  \item Quand $\psi$ va de $0$ à $-\pi$ (sens horaire), le point parcourt le \textbf{demi-cercle du bas}.
        Il est strictement en dessous, donc $\sin\psi<0$.
  \item Aux extrémités $\psi=0$ et $\psi=\pm\pi$, le point est sur l'axe horizontal : $\sin\psi=0$.
\end{itemize}

\subsection{Démonstration analytique}

\begin{lemme}\label{lem:sinpos}
Pour $0<\psi<\pi$, on a $\sin\psi>0$.
\end{lemme}
\begin{proof}
\textbf{Étape 2.a : sur $\left]0,\frac{\pi}{2}\right]$.}
La dérivée de $\sin$ est $\cos$. Pour $0\le x<\frac\pi2$, on a $\cos x>0$, donc $\sin$ est strictement
croissante sur $\left[0,\frac\pi2\right]$. Comme $\sin 0 = 0$ :
\begin{equation}
0<x\le\tfrac{\pi}{2}\quad\Longrightarrow\quad \sin x>\sin 0 = 0 .
\end{equation}

\textbf{Étape 2.b : sur $\left]\frac{\pi}{2},\pi\right[$.}
On utilise la symétrie $\sin(\pi-x) = \sin x$, qui découle de la formule
$\sin(a-b)=\sin a\cos b-\cos a\sin b$ :
\begin{equation}
\sin(\pi-x) = \sin\pi\cos x - \cos\pi\sin x = 0\cdot\cos x-(-1)\sin x = \sin x .
\end{equation}
Si $\frac\pi2<x<\pi$, alors $0<\pi-x<\frac\pi2$. L'étape 2.a donne $\sin(\pi-x)>0$, donc $\sin x>0$.

\textbf{Conclusion.} En réunissant 2.a et 2.b : $\sin\psi>0$ pour tout $0<\psi<\pi$.
\end{proof}

%=====================================================================
\section{Étape 3 : $\sin\psi<0$ pour $-\pi<\psi<0$}
%=====================================================================

\begin{lemme}[Imparité du sinus]\label{lem:impair}
Pour tout $x$ : $\sin(-x) = -\sin x$.
\end{lemme}
\begin{proof}
Sur le cercle trigonométrique, l'angle $-x$ est le symétrique de l'angle $x$ par rapport à l'axe horizontal.
L'abscisse est conservée et l'ordonnée change de signe.
(Ou bien : $\sin$ est la fonction impaire $x - \frac{x^3}{3!} + \frac{x^5}{5!}-\dots$, qui ne contient que des puissances impaires.)
\end{proof}

\textbf{Étape 3.a.} Soit $-\pi<\psi<0$. On pose $u = -\psi$. Alors $0<u<\pi$.

\textbf{Étape 3.b.} D'après le lemme~\ref{lem:sinpos} : $\sin u>0$.

\textbf{Étape 3.c.} D'après le lemme~\ref{lem:impair} :
\begin{equation}
\sin\psi = \sin(-u) = -\sin u <0 .
\end{equation}

%=====================================================================
\section{Conclusion de la démonstration}
%=====================================================================

On rassemble les trois étapes dans un tableau :

\begin{center}
\renewcommand{\arraystretch}{1.4}
\begin{tabular}{|c|c|c|c|c|}
\hline
Intervalle & $\sign(\psi)$ & $\sin\psi$ & $\sign(\sin\psi)$ & $\sign(\varepsilon)$ \\
\hline
$-\pi<\psi<0$ & $-1$ & $<0$ (étape 3) & $-1$ & $-1$ \\
$\psi=0$ & $0$ & $=0$ & $0$ & $0$ \\
$0<\psi<\pi$ & $+1$ & $>0$ (étape 2) & $+1$ & $+1$ \\
\hline
\end{tabular}
\end{center}

La deuxième et la dernière colonnes sont identiques. Avec \eqref{eq:etape1} :

\begin{resultat}
\begin{equation}
\sign(\varepsilon) = \sign(\sin\psi) = \sign(\psi)\qquad\text{pour } |\psi|<\pi .
\end{equation}
(L'égalité est aussi vraie en $\psi=0$, où les deux membres valent $0$.)
\end{resultat}

\subsection*{Illustration graphique}

\begin{center}
\begin{tikzpicture}
\begin{axis}[
  width=0.95\textwidth, height=7cm,
  xmin=-2*pi-0.2, xmax=2*pi+0.2, ymin=-1.45, ymax=1.45,
  axis lines=middle, axis on top, xlabel={$\psi$}, ylabel={$\varepsilon/V_m=\sin\psi$},
  xlabel style={at={(1,0.5)},anchor=west},
  xtick={-6.2832,-4.7124,-3.1416,-1.5708,0,1.5708,3.1416,4.7124,6.2832},
  xticklabels={$-2\pi$,$-\frac{3\pi}{2}$,$-\pi$,$-\frac{\pi}{2}$,$0$,$\frac{\pi}{2}$,$\pi$,$\frac{3\pi}{2}$,$2\pi$},
  ytick={-1,1}, samples=300, domain=-2*pi:2*pi,
]
  \fill[green!12] (axis cs:0,-1.45) rectangle (axis cs:pi,1.45);
  \fill[red!10] (axis cs:-pi,-1.45) rectangle (axis cs:0,1.45);
  \fill[gray!15] (axis cs:pi,-1.45) rectangle (axis cs:2*pi+0.2,1.45);
  \fill[gray!15] (axis cs:-2*pi-0.2,-1.45) rectangle (axis cs:-pi,1.45);
  \addplot[very thick,blue] {sin(deg(x))};
  \node[green!35!black,align=center,font=\small] at (axis cs:pi/2,-0.6) {$\psi>0$ et $\varepsilon>0$\\ \textbf{même signe}};
  \node[red!60!black,align=center,font=\small] at (axis cs:-pi/2,0.6) {$\psi<0$ et $\varepsilon<0$\\ \textbf{même signe}};
  \node[black!70,align=center,font=\small] at (axis cs:3*pi/2,0.6) {$\psi>0$ mais $\varepsilon<0$\\ \textbf{signes opposés}};
  \node[black!70,align=center,font=\small] at (axis cs:-3*pi/2,-0.6) {$\psi<0$ mais $\varepsilon>0$\\ \textbf{signes opposés}};
  \addplot[only marks,mark=*,mark size=2.5pt,green!50!black] coordinates {(0,0)};
  \addplot[only marks,mark=o,mark size=3pt,red,thick] coordinates {(-pi,0) (pi,0)};
\end{axis}
\end{tikzpicture}
\end{center}
Dans la zone colorée $|\psi|<\pi$ (vert et rouge), $\varepsilon$ et $\psi$ ont toujours le même signe.
Dans les zones grises ($|\psi|>\pi$), ce n'est plus vrai.

%=====================================================================
\section{Pourquoi la condition $|\psi|<\pi$ est indispensable}
%=====================================================================

\subsection{Contre-exemples}
\begin{itemize}
  \item $\psi = \frac{3\pi}{2}$ : $\psi>0$ mais $\sin\frac{3\pi}{2} = -1<0$. Les signes sont \textbf{opposés}.
  \item $\psi = \pi$ : $\psi>0$ mais $\sin\pi = 0$, donc $\sign(\varepsilon) = 0\neq+1$.
  \item $\psi = 2\pi$ : $\psi>0$ mais $\sin 2\pi = 0$.
\end{itemize}

\subsection{Explication}
Le sinus est $2\pi$-périodique : le détecteur ne \og voit\fg{} pas la différence entre $\psi$ et $\psi+2\pi$.
Physiquement, une erreur de $\frac{3\pi}{2}$ (retard de $270^\circ$) est la même chose qu'une erreur de
$-\frac{\pi}{2}$ (avance de $90^\circ$). Le détecteur choisit toujours l'interprétation la plus courte.

\subsection{Conséquence pratique}
On ramène toujours l'erreur dans l'intervalle $]-\pi,\pi]$ :
\begin{equation}
\psi_{\text{ramené}} = \psi - 2\pi\left\lfloor\frac{\psi+\pi}{2\pi}\right\rfloor\in\,]-\pi,\pi] .
\end{equation}
Par périodicité, $\sin\psi_{\text{ramené}} = \sin\psi$. On a donc toujours
\begin{equation}
\sign(\varepsilon) = \sign(\psi_{\text{ramené}}).
\end{equation}
C'est pour la même raison que l'on calcule $\hth$ modulo $2\pi$ dans l'implémentation.

%=====================================================================
\section{Signification physique et rôle dans la PLL}
%=====================================================================

\begin{center}
\renewcommand{\arraystretch}{1.4}
\begin{tabular}{|c|c|c|c|}
\hline
$\psi=\theta-\hth$ & Situation & $\sign(\varepsilon)$ & Action de la PLL \\
\hline
$>0$ & $\hth<\theta$ : estimation \textbf{en retard} & $+1$ & $\dot z = +k_2$, $\hat\omega$ augmente : $\hth$ accélère\\
$<0$ & $\hth>\theta$ : estimation \textbf{en avance} & $-1$ & $\dot z = -k_2$, $\hat\omega$ diminue : $\hth$ ralentit\\
$=0$ & verrouillée & $0$ & aucune correction\\
\hline
\end{tabular}
\end{center}

\begin{idee}
La PLL ne connaît pas $\psi$ : elle ne mesure que $\varepsilon$. Cette propriété lui garantit que
\textbf{le signe de ce qu'elle mesure ($\varepsilon$) est le signe de ce qu'elle doit corriger ($\psi$)}.
La correction se fait donc toujours dans le bon sens.
\end{idee}

\subsection*{Utilisation dans la preuve de stabilité}
Dans les équations de l'erreur, on remplace
\begin{equation}
\dot z = k_2\sign(e) = k_2\sign(\sin x_1) = k_2\sign(x_1),
\end{equation}
avec $e = \varepsilon/\hat V_m$ et $x_1=\psi$. On obtient ainsi
\begin{equation}
\dot x_2 = -k_2\sign(x_1)+\dot\omega ,
\end{equation}
qui est exactement la forme canonique du Super-Twisting. Sans cette propriété, $\sign(e)$ ne pourrait pas
s'écrire en fonction de $x_1$ seul, et l'analyse de Lyapunov ne pourrait pas être menée.

\end{document}
